\documentclass[12pt]{article}
\usepackage[utf8]{inputenc}
\usepackage{amsmath,amssymb,amsthm}
\usepackage{algorithm}
\usepackage{algpseudocode}
\usepackage{geometry}
\geometry{a4paper, margin=1in}

\newtheorem{theorem}{Theorem}
\newtheorem{lemma}[theorem]{Lemma}
\newtheorem{corollary}[theorem]{Corollary}
\newtheorem{proposition}[theorem]{Proposition}
\theoremstyle{definition}
\newtheorem{definition}[theorem]{Definition}

\title{Problem 38: Pandigital Multiples}
\author{Project Euler Solution}
\date{}

\begin{document}
\maketitle

\section{Problem Statement}
For a fixed positive integer $x$, form the concatenated product of $x$ with $(1, 2, \ldots, n)$ for some $n > 1$. What is the largest 1-to-9 pandigital 9-digit number achievable?

\section{Definitions}

\begin{definition}
The \emph{concatenated product} of $x$ with $(1, 2, \ldots, n)$ is
\[
    \mathrm{CP}(x, n) = \overline{(x \cdot 1) \| (x \cdot 2) \| \cdots \| (x \cdot n)},
\]
where $\|$ denotes decimal string concatenation.
\end{definition}

\begin{definition}
An integer is \emph{1-to-9 pandigital} if its decimal digits are exactly $\{1,2,3,4,5,6,7,8,9\}$, each appearing once.
\end{definition}

\section{Theoretical Development}

\begin{theorem}[Digit count constraint]
\label{thm:digits}
For $\mathrm{CP}(x,n)$ to be 9-digit pandigital with $n \ge 2$, we need $\sum_{k=1}^{n} (\lfloor \log_{10}(kx) \rfloor + 1) = 9$.
\end{theorem}
\begin{proof}
The concatenation has $\sum_{k=1}^n d(kx)$ digits, which must equal~9 for pandigitality.
\end{proof}

\begin{theorem}[Feasible parameter space]
\label{thm:space}
The constraint restricts $(x,n)$ to finitely many families:
\begin{center}
\begin{tabular}{ccc}
\hline
$n$ & Digit pattern & Range of $x$ \\
\hline
2 & $(4,5)$ & $5000 \le x \le 9999$ \\
3 & $(3,3,3)$ & $100 \le x \le 333$ \\
4 & $(2,2,2,3)$ & $25 \le x \le 33$ \\
5 & $(1,2,2,2,2)$ & $5 \le x \le 9$ \\
9 & $(1,1,\ldots,1)$ & $x = 1$ \\
\hline
\end{tabular}
\end{center}
\end{theorem}
\begin{proof}
For $n=2$: $d(x) + d(2x) = 9$. With $d(x)=4$, we need $d(2x)=5$, i.e., $2x \ge 10000$ so $x \ge 5000$. For $d(x) \ge 5$: total $\ge 10 > 9$. The other cases follow by analogous arithmetic. For $n \ge 10$, even $x=1$ yields $\ge 11$ digits.
\end{proof}

\begin{lemma}[Leading-digit maximization]
\label{lem:leading}
Since $\mathrm{CP}(x,n)$ begins with $x$, maximizing the result requires $x$ to begin with the digit~$9$.
\end{lemma}
\begin{proof}
For two 9-digit numbers, the one with the lexicographically larger prefix is numerically larger. The first $d(x)$ digits of $\mathrm{CP}(x,n)$ equal~$x$.
\end{proof}

\begin{theorem}[Optimality]
\label{thm:optimal}
The maximum pandigital concatenated product is $\mathrm{CP}(9327, 2) = 932718654$.
\end{theorem}
\begin{proof}
By Lemma~\ref{lem:leading} and Theorem~\ref{thm:space}, the $n=2$ family with $x \in [9000, 9999]$ dominates all other families (the best $n=5$ case gives $918273645 < 932\ldots$). At $x=9327$:
\[
    \mathrm{CP}(9327, 2) = \overline{9327 \| 18654} = 932718654.
\]
The digit multiset $\{9,3,2,7,1,8,6,5,4\} = \{1,\ldots,9\}$ confirms pandigitality. Exhaustive search over $x \in [9328, 9999]$ yields no pandigital result, and no other family produces a larger value.
\end{proof}

\section{Algorithm}

\begin{algorithmic}[1]
\Function{LargestPandigitalMultiple}{}
    \State $\mathit{best} \gets 0$
    \For{$x = 1$ to $9999$}
        \State $s \gets \text{``''}$, $n \gets 0$
        \While{$|s| < 9$}
            \State $n \gets n + 1$
            \State $s \gets s \| \mathrm{str}(x \cdot n)$
        \EndWhile
        \If{$n \ge 2$ \textbf{and} $|s| = 9$ \textbf{and} $s$ is pandigital}
            \State $\mathit{best} \gets \max(\mathit{best},\; \mathrm{int}(s))$
        \EndIf
    \EndFor
    \State \Return $\mathit{best}$
\EndFunction
\end{algorithmic}

\section{Complexity}

\begin{proposition}
The algorithm runs in $O(9999 \times 9) = O(9 \times 10^4)$ time and $O(1)$ auxiliary space.
\end{proposition}
\begin{proof}
The outer loop runs $9999$ times. Each inner loop performs at most $9$ concatenations. The pandigital check on a $9$-character string costs $O(1)$. No auxiliary data structures are needed.
\end{proof}

\section{Result}
\[
    \boxed{932718654}
\]

\section{Answer}

\[
\boxed{932718654}
\]

\end{document}
